\documentclass[11pt]{article}
\usepackage[margin=1in]{geometry}
\usepackage{amsmath,amssymb}
\usepackage{booktabs}
\usepackage{enumitem}
\usepackage{array}
\usepackage{hyperref}
\usepackage{parskip}
\setlist{nosep,leftmargin=*}

\begin{document}

\noindent{\Large\bfseries Anomaly Detection for Anti-Money Laundering: Literature Review}\\[2pt]
{\itshape Internal research note -- January 2026 (v1)}

\medskip

\textbf{Scope.} Survey of ML/DL methods for AML anomaly detection, focusing on 2020--2026. Covers graph-based, deep learning, traditional ML, contextual/KYC, feature engineering, datasets, and emerging paradigms. Not a formal paper---a working reference for project planning.

\textbf{Audience.} Data scientist with anomaly detection experience starting an AML project.

\bigskip
\noindent{\large\bfseries Table of Contents}
\medskip
\setcounter{tocdepth}{1}
\makeatletter
\@starttoc{toc}
\makeatother

\newpage

%==============================================================================
\section{Graph and Network-Based Approaches}
%==============================================================================

Financial transactions form natural graphs (nodes = accounts, edges = transfers). Graph methods detect laundering patterns invisible in tabular data.

\subsection{GNNs for AML}

The Elliptic dataset (Weber et al., 2019, KDD Workshop) is the foundational benchmark: 203K Bitcoin transactions, 234K edges, $\sim$2\% illicit. GCN achieved $\sim$0.97 AUC, showing graph structure adds value over node features alone.

Cheng et al.\ (2024, \textit{Frontiers of Computer Science}, arXiv:2411.05815) survey 100+ GNN methods for financial fraud, finding GNNs substantially outperform traditional approaches.

\begin{center}
\small
\begin{tabular}{llccl}
\toprule
\textbf{Method} & \textbf{Dataset} & \textbf{AUC} & \textbf{F1 (illicit)} & \textbf{Note} \\
\midrule
Random Forest & Elliptic & 0.96 & 0.47 & Node features only \\
GCN & Elliptic & 0.97 & 0.52 & +graph structure \\
GAT & Elliptic & 0.97 & 0.51 & Attention-based \\
GraphSAGE & Elliptic & 0.97 & 0.54 & Inductive \\
EvolveGCN-H & Elliptic & 0.97 & 0.58 & +temporal \\
ATGAT & Elliptic++ & 0.91 & -- & 9.2\% over XGBoost \\
CARE-GNN & YelpChi & 0.94 & 0.86 & Anti-camouflage \\
PC-GNN & YelpChi & 0.97 & 0.88 & Imbalanced learning \\
GraphSAGE & AMLSim & 0.85--0.95 & 0.70--0.85 & Varies by pattern \\
\bottomrule
\end{tabular}
\end{center}

Key architectures:
\begin{itemize}
\item \textbf{GCN/GAT/GraphSAGE}: Standard baselines. GraphSAGE preferred for inductive (unseen node) settings. Recent temporal-aware GAT (ATGAT) on Elliptic++ achieves 0.913 AUC.
\item \textbf{Improved HybridNet} (2025): Combines GraphSAGE + GATv2 + temporal encoding.
\item \textbf{LineMVGNN} (2025): Line-graph multi-view GNN for directed financial graphs.
\item \textbf{LayerWeighted-GCN}: Adaptive layer weighting across GNN depths.
\end{itemize}

\subsection{Self-Supervised and Contrastive Learning on Graphs}

Label scarcity is fundamental in AML (labeled SAR data is rare and sensitive):

\begin{itemize}
\item \textbf{LaundroGraph} (Cardoso et al., 2022, ECML-PKDD Workshop): Self-supervised graph embeddings $\to$ XGBoost. F1 improved $\sim$5--10\% over raw tabular features.
\item \textbf{Bit-CHetG} (2024): Subgraph contrastive learning on heterogeneous graphs for laundering group detection.
\item \textbf{HCLNet} (2025): Hypergraph contrastive learning for group-level fraud semantics.
\item \textbf{Causal-DHG} (2024): Dynamic hypergraph + causal intervention + multi-view contrastive learning.
\item \textbf{GRAD} (2025, TheWebConf): Guided relation diffusion + supervised graph contrastive learning from limited labels.
\item \textbf{SpaceGNN} (2025, ICLR): Multi-space projections with distance-aware propagation.
\end{itemize}

\subsection{Heterogeneous Graphs}

Real financial networks have multiple entity types (persons, accounts, merchants, devices) and relation types:

\begin{itemize}
\item \textbf{CARE-GNN} (Dou et al., 2020, CIKM): Camouflage-resistant; RL-based relation selector. AUC 0.94 on YelpChi.
\item \textbf{PC-GNN} (Liu et al., 2021, WWW): Pick-and-choose sampling for imbalance. AUC 0.97.
\item \textbf{xFraud} (Rao et al., 2022, VLDB): Industrial-scale heterogeneous GNN deployed at eBay scale with explainability. Billions of nodes.
\item \textbf{HAN}: Hierarchical attention (node-level + meta-path semantic-level). 2--8\% AUC improvement over homogeneous GNNs.
\item \textbf{HTGNN} (2025): Heterogeneous temporal GNN integrating spatial, temporal, and semantic information for real-time detection.
\end{itemize}

\subsection{Temporal and Dynamic Graphs}

\begin{itemize}
\item \textbf{EvolveGCN} (Pareja et al., 2020, AAAI): RNN evolves GCN weights over time. F1 0.58 vs.\ 0.52 for static GCN on Elliptic.
\item \textbf{TGN} (Rossi et al., 2020): Memory-augmented temporal graph network with per-node memory. Adapted for AML by subsequent works.
\item \textbf{GTCT} (2025): Graph-Temporal Contrastive Transformer for transactional behavior modeling.
\item \textbf{Continual Graph Learning for AML} (Deprez et al., 2025, \textit{WIREs Computational Statistics}): Reviews replay-based, regularization-based, and architecture-based strategies against catastrophic forgetting.
\item Streaming GNNs from Ant Group, PayPal for real-time processing on evolving graphs.
\end{itemize}

\subsection{Subgraph and Community Detection}

Money laundering involves coordinated account networks. Node-level detection misses this:

\begin{itemize}
\item \textbf{FlowScope} (Li et al., 2020, KDD): Dense subgraph detection for scatter-gather laundering patterns. Orders of magnitude faster than subgraph enumeration.
\item \textbf{SMoTeF} (2024, \textit{Applied Intelligence}): Smurf detection using temporal order and flow analysis. Applied to 180M+ transactions, 31M accounts.
\item \textbf{GARG-AML} (2025): Scalable, interpretable smurfing detection. Outperforms SOTA on large-scale networks.
\item \textbf{AMLGaurd} (2025): Unsupervised GAE-based detection via reconstruction discrepancy.
\item \textbf{Graph Feature Preprocessor} (2024, KDD): Real-time subgraph pattern mining on streaming transactions. Multicore parallelism.
\item \textbf{The Shape of Money Laundering} (2024): Subgraph representation learning on Elliptic2 revealing laundering scheme shapes (peeling chains, nested services).
\end{itemize}

%==============================================================================
\section{Deep Learning (Non-Graph)}
%==============================================================================

\subsection{Sequence Models (LSTM, GRU)}

Transaction sequences are temporal $\Rightarrow$ RNNs are natural:

\begin{itemize}
\item LSTM/GRU outperform SimpleRNN (vanishing gradients). GRU often preferred for AML: fewer params, AUC-ROC $>$ 0.7. GRU matches or beats LSTM on most AML benchmarks.
\item Attention-enhanced LSTM (Niu et al., 2021): Attention focuses on most suspicious subsequences.
\item Jensen and Iosifidis (2023): Replace rule-based triggers with auto-extracted sequential features at bank level.
\end{itemize}

\subsection{Transformers}

\begin{itemize}
\item \textbf{Tang and Liu (2024)}: Unified Transformer-BDI with distributed knowledge distillation. F1 92.87\%, AUC 96.73\%.
\item \textbf{TREASURE} (KDD 2026): Transformer foundation model for high-volume transaction understanding.
\item \textbf{TransactionGPT} (2025): GPT-style pretraining on payment data for behavioral anomaly detection.
\item \textbf{FraudGT} (2024, ICAIF): Graph Transformer for financial fraud.
\item \textbf{TabTransformer, TransTab, BERT, RoBERTa} applied to structured transaction data (2024).
\item \textbf{Spatial-Temporal-Aware Graph Transformer} (2024, IEEE TII).
\end{itemize}

\subsection{Autoencoders and VAEs}

Trained on normal transactions, high reconstruction error = anomaly:

\begin{itemize}
\item \textbf{VAE for CC Anomaly} (2024, \textit{Big Data Mining and Analytics}): CNN-based VAE; hybrid VAE + classifier models.
\item \textbf{GAN+VAE Hybrid} (2025): Combines GANs and VAEs for large payment flow anomaly detection.
\item \textbf{LVAESC-SFD} (2026): Correlation-based feature selection + VAE, 98.24\% accuracy.
\item \textbf{Denoising Autoencoder} (Bundesbank, 2023): Detection + interpretability on financial tabular data.
\item \textbf{Self-Attention + VAE} (2024, \textit{Expert Systems with Applications}): Fund transfer fraud with attention for long-range dependencies.
\item Shallow autoencoders (1--2 layers) often sufficient for tabular AML. AUC $\sim$0.85--0.90 on CC fraud, $\sim$0.80--0.85 on AML-specific data.
\end{itemize}

\subsection{GANs for AML}

Dual purpose: synthetic data generation and direct anomaly detection.

\begin{itemize}
\item \textbf{WGAN for data augmentation}: Synthetic fraudulent transactions improve F1 by 10--15\%.
\item \textbf{SMOTE-GAN hybrids} (2023): SMOTified-GAN, SMOTE+GAN, GANified-SMOTE variants.
\item \textbf{Advanced R-GAN} (2024): Regularized GAN with spectral normalization against mode collapse.
\item \textbf{CTGAN, Copula GAN, TVAE}: State-of-the-art synthetic data tools. 96--99\% utility equivalence to production data.
\item Alan Turing Institute project (commenced autumn 2024): GAN-based synthetic AML data for UK FCA initiative.
\end{itemize}

\subsection{LLMs and Foundation Models}

Newest frontier (2025--2026):

\begin{itemize}
\item \textbf{FLAG} (KDD 2025): LLM-enhanced GNN for fraud detection. Uses large language models to generate semantic node features from raw transaction metadata, which are then fused with structural graph embeddings. Demonstrates that LLM-derived features can substitute for manual feature engineering in graph-based AML pipelines.
\item \textbf{GuARD} (KDD 2025): Text-rich, graph-informed language model for anomaly detection. Jointly encodes unstructured text fields (e.g., transaction memos, KYC notes) and graph topology into a unified representation. Particularly relevant for AML settings where free-text fields carry signals that tabular features miss.
\item \textbf{DGP} (AAAI 2026): Dual-granularity prompting with graph-enhanced LLMs. Introduces two levels of prompting---transaction-level and subgraph-level---to capture both local and structural anomaly signals. Achieves competitive performance with GNNs while offering natural-language explanations as a by-product.
\item \textbf{AuditCopilot} (NeurIPS 2025 Workshop): LLMs for bookkeeping fraud. Applies instruction-tuned LLMs to detect anomalies in accounting records and journal entries. Highlights the potential for LLMs in adjacent financial crime domains beyond transaction monitoring.
\item \textbf{Co-Investigator AI} (2025): Agentic framework for faster, more accurate SAR generation. Combines retrieval-augmented generation with structured case data to draft SAR narratives, reducing analyst writing time. Positions LLMs as analyst-assistance tools rather than standalone detectors.
\item \textbf{FinGPT} (IJCAI 2023): Open-source financial LLM. Pre-trained on financial corpora and fine-tunable for downstream tasks including sentiment analysis, risk assessment, and fraud classification. Provides a foundation model baseline for financial NLP tasks, though AML-specific fine-tuning remains underexplored.
\end{itemize}

%==============================================================================
\section{Traditional ML}
%==============================================================================

\subsection{Gradient Boosting}

Dominates tabular AML. Consistent with Grinsztajn et al.\ (2022, NeurIPS): tree-based models outperform DL on structured data.

\begin{center}
\small
\begin{tabular}{llll}
\toprule
\textbf{Method} & \textbf{Typical AUC} & \textbf{Speed} & \textbf{Key Advantage} \\
\midrule
XGBoost & 0.95--0.99 & Medium & Best overall tabular; built-in feature importance \\
LightGBM & 0.95--0.98 & Fastest & 3--5$\times$ faster; high-cardinality cats \\
CatBoost & 0.95--0.98 & Slow & Native categoricals; ordered boosting \\
Stacking & 0.96--0.99 & Slowest & Best combined performance \\
\bottomrule
\end{tabular}
\end{center}

\subsection{Unsupervised Methods}

\begin{center}
\small
\begin{tabular}{lcll}
\toprule
\textbf{Method} & \textbf{AUC} & \textbf{Strengths} & \textbf{Weaknesses} \\
\midrule
Isolation Forest & 0.78--0.85 & No labels; fast; scalable & Lower precision standalone \\
Extended IF & +1--2\% & Better correlated features & -- \\
LOF & 0.72--0.80 & Local anomalies & Doesn't scale; k-sensitive \\
One-Class SVM & 0.70--0.80 & Only needs normal class & Kernel choice; doesn't scale \\
DBSCAN & 0.60--0.75 & Structural outliers & Low precision \\
PCA & 0.70--0.85 & Fast; interpretable & Misses nonlinear \\
\bottomrule
\end{tabular}
\end{center}

\subsection{Hybrid Rule + ML}

Industry standard. Pure ML replacement of rules is rare due to regulatory requirements:

\begin{itemize}
\item \textbf{Three-tier FRAML architecture} (2023): (1) hard rules for regulatory compliance, (2) soft rules for known typologies, (3) ML for emerging patterns.
\item Rocha-Salazar et al.\ (2021): Rules generate alerts; RF re-scores. Alert volume $-70\%$; SAR precision from 2\% (rules) to 15\% (hybrid).
\item Soltani Delgosha and Hajiheydari (2021): Staged rule+ML: $-40\%$ false positives while maintaining recall.
\item Regulators mandate explainable triggers $\Rightarrow$ rules for audit trail, ML for ranking.
\end{itemize}

%==============================================================================
\section{Contextual Anomaly Detection}
%==============================================================================

\textbf{Core idea.} Contextual anomaly detection models $P(\mathbf{y}|\mathbf{c})$---expected behavior $\mathbf{y}$ given context $\mathbf{c}$---rather than a global $P(\mathbf{x})$. A \$50K wire is normal for a corporate treasury but anomalous for a student account.

Contextual methods can be broadly categorized by how they incorporate context:

\begin{center}
\small
\begin{tabular}{p{3.2cm}p{4.5cm}p{4cm}}
\toprule
\textbf{Category} & \textbf{Approach} & \textbf{Examples} \\
\midrule
Peer-normalized & Normalize features by peer group statistics before applying a standard detector & Peer Z-score + IF, Bolton \& Hand (2001) \\
Hard contextual & Partition data by discrete context (e.g., customer segment), fit separate models per group & Segment-level IF, per-group thresholds \\
Soft contextual & Model $P(\mathbf{y}|\mathbf{c})$ directly with continuous context conditioning & CAD, ROCOD, QCAD, CWAE \\
Learned contextual & Automatically discover which features are context vs.\ behavioral & ConQuest, CWAE bilevel optimization \\
Uncertainty-aware & Provide calibrated confidence on anomaly scores given context density & GP-based contextual AD (2025) \\
\bottomrule
\end{tabular}
\end{center}

\subsection{Foundational Methods}

\begin{center}
\small
\begin{tabular}{p{2cm}lp{4cm}lc}
\toprule
\textbf{Method} & \textbf{Year} & \textbf{Key Idea} & \textbf{Venue} & \textbf{Avg AUC} \\
\midrule
CAD & 2007 & Conditional Gaussian mixture; anomaly = low $P(\mathbf{y}|\mathbf{c})$ & IEEE TKDE & -- \\
ROCOD & 2016 & Local + global models; adaptive weighting by context density & CIKM & 0.75--0.82 \\
QCAD & 2023 & Quantile regression forests; explainable via quantile intervals & DAMI & 0.78--0.85 \\
ConQuest & 2024 & Auto-discovers contexts via genetic algorithm; MCAF scoring & IJDSA 2025 & SOTA on 25 datasets \\
\bottomrule
\end{tabular}
\end{center}

\subsection{Deep Contextual Methods}

\subsubsection{Conditional Wasserstein Autoencoder (CWAE)}

\textbf{King et al.\ (2025), ``Contextual Learning for Anomaly Detection in Tabular Data''} (arXiv:2509.09030, submitted to TMLR). This is the most directly relevant recent paper for deep contextual AD on tabular data.

\textbf{Architecture.} Splits features into context $\mathbf{c}$ and content/behavioral $\mathbf{y}$. Deterministic encoder $\mathbf{z} = E_\phi([\mathbf{y}; \mathbf{c}])$, decoder $\hat{\mathbf{y}} = D_\theta([\mathbf{z}; \mathbf{c}])$ reconstructs behavioral features only. Loss:
\begin{equation*}
\mathcal{L} = \sum_{i=1}^{d} \mathrm{CE}(y_i, \hat{y}_i) + \lambda \cdot \mathrm{MMD}^2(P(\mathbf{z}), Q(\mathbf{z}|\mathbf{x}))
\end{equation*}
Uses MMD with multi-scale IMQ kernel instead of KL divergence $\Rightarrow$ deterministic encoder, no sampling noise, bounded gradients.

\textbf{Context selection.} Bilevel optimization: outer loop iterates over candidate context features; inner loop trains CWAE for 1 epoch; selects context minimizing joint validation loss $\mathcal{L}_{\mathrm{val}} = -\log P(\mathbf{y}, \mathbf{c})$. Leverages the finding that early-trained models proxy generalization quality.

\textbf{Theoretical justification.} Variance decomposition: $\mathrm{Var}(\mathbf{y}) = \mathbb{E}_{\mathbf{c}}[\mathrm{Var}(\mathbf{y}|\mathbf{c})] + \mathrm{Var}_{\mathbf{c}}(\mathbb{E}[\mathbf{y}|\mathbf{c}])$. A context-conditional detector only needs to model $\mathrm{Var}(\mathbf{y}|\mathbf{c})$, which is typically much smaller than total $\mathrm{Var}(\mathbf{y})$ $\Rightarrow$ higher SNR for anomaly detection.

\textbf{Results} on 8 benchmarks (Bank, Beth, Census, CMC, KDD, LANL, SF, Spotify):

\begin{center}
\small
\begin{tabular}{lccc}
\toprule
\textbf{Method} & \textbf{Avg AUC} & \textbf{Avg Rank} & \textbf{Type} \\
\midrule
CWAE (proposed) & \textbf{0.797} & \textbf{2.75} & Contextual deep \\
WAE (non-contextual) & 0.760 & 4.25 & Global deep \\
CVAE & 0.745$^*$ & -- & Contextual deep \\
CMDN & 0.756$^*$ & -- & Contextual deep \\
CADI & 0.734$^*$ & -- & Contextual tree \\
RCA & 0.769 & 3.38 & Global deep \\
DIF & 0.758 & 3.50 & Global tree \\
SLAD & 0.695 & 5.88 & Global deep \\
\bottomrule
\multicolumn{4}{l}{\scriptsize $^*$Estimated from ablation tables; exact ordering varies by dataset.}
\end{tabular}
\end{center}

Key finding: with oracle context selection, CWAE achieves avg AUC 0.839 (rank 1 on 6/8 datasets). The gap between automatic and oracle context selection highlights that \textit{context feature choice} is the critical design decision.

\textbf{CWAE vs.\ CVAE comparison.} CWAE outperforms CVAE (+11.69\% avg improvement over non-contextual baseline vs.\ CVAE's +6.21\%). CWAE's deterministic encoder avoids CVAE's numerical instability: CVAE requires $\log\sigma^2$ clamping to prevent divergence on skewed financial data, and Monte Carlo sampling adds inference noise.

\subsubsection{Conditional VAE (CVAE) for Anomaly Detection}

CVAE models $q_\phi(\mathbf{z}|\mathbf{x}, \mathbf{c})$ with KL regularization. Applied to contextual AD by conditioning the generative model on context features. Three instability sources on financial data: (1) $\log\sigma^2 \to -\infty$ when encoder predicts near-zero variance, (2) $\log\sigma^2 \to +\infty$ causes overflow, (3) reparameterization amplifies noise when $\sigma$ is extreme. The KL term also forces a ``fuzzy'' Gaussian latent space that may reconstruct anomalies too well.

\textbf{Disentangled CVAE} (IEEE 2024): Combines $\beta$-VAE, CVAE, and total correlation to disentangle context from anomaly signals. Addresses information loss while incorporating known sources of variation.

\subsubsection{Other Deep Contextual Methods}

\begin{itemize}
\item \textbf{CMDN} (Conditional Mixture Density Network): Neural network-parameterized Gaussian mixture for $P(\mathbf{y}|\mathbf{c})$. More flexible than single Gaussian but harder to train. +7.47\% improvement over non-contextual (King et al., 2025).
\item \textbf{CADI} (Contextual Anomaly Detection using Isolation Forest): Extends IF with density-aware splits conditioned on context. +4.18\% improvement. Tree-based $\Rightarrow$ more interpretable.
\end{itemize}

\subsection{Peer Group and Peer-Normalized Methods}

The operational form of contextual AD in AML: compare each entity to its ``peers'' rather than the global population.

\textbf{Bolton and Hand (2001).} Foundational peer group analysis for fraud detection. Group entities by context features, flag deviations from peer behavior. The core idea: compute per-entity statistics relative to peer group mean and standard deviation.

\textbf{Peer-normalized Z-score approach.} For each entity $i$ with context $\mathbf{c}_i$ and behavioral feature $y_i$:
\begin{equation*}
z_i = \frac{y_i - \mu_{\mathrm{peer}}(\mathbf{c}_i)}{\sigma_{\mathrm{peer}}(\mathbf{c}_i)}
\end{equation*}
where $\mu_{\mathrm{peer}}$ and $\sigma_{\mathrm{peer}}$ are estimated from context-similar entities (KNN with kernel weighting, or hard peer group assignment). This is effectively the simplest form of contextual AD.

\textbf{Peer-Normalized Isolation Forest.} Apply IF to peer-normalized features rather than raw features. This operationalizes contextual AD without changing the detector---only the feature representation changes. The normalization removes context-induced location/scale shifts, exposing within-context anomalies. Computationally cheap: no model retraining needed per context.

\textbf{Uncertainty-Aware Contextual AD} (arXiv:2507.04490, 2025). Extends the peer-normalized Z-score by treating it as a \textit{random variable} with confidence intervals. Uses two heteroscedastic Gaussian process regression models---one for the conditional mean $\mathbb{E}[y|\mathbf{c}]$, one for the conditional variance $\mathrm{Var}(y|\mathbf{c})$. Provides calibrated uncertainty: wide confidence intervals $\Rightarrow$ low reliability of anomaly assessment. Outperforms QCAD and ROCOD on benchmark datasets and a real-world cardiology application. Key advantage: enables adaptive, uncertainty-driven decision-making rather than hard thresholds.

\subsection{When Does Context Help?}

Not always. Contextual methods add value specifically when anomalies are \textit{contextually unusual but globally normal}:

\begin{center}
\small
\begin{tabular}{p{3cm}p{4.5cm}p{4cm}}
\toprule
\textbf{Anomaly Type} & \textbf{Example} & \textbf{Best Detector} \\
\midrule
Global & \$10M single transfer (unusual for anyone) & Standard IF, XGBoost \\
Contextual & \$50K wire from a student account (normal for corp.) & Peer-normalized IF, CWAE, QCAD \\
Structural & Fan-out $\to$ fan-in chain (smurfing) & GNN, FlowScope \\
Temporal & Dormant account suddenly active & LSTM, TGN \\
\bottomrule
\end{tabular}
\end{center}

Key empirical finding: the magnitude of improvement from contextual methods depends on (1) the \textit{type} of anomaly present and (2) the contamination rate. Under context-mismatch contamination, contextual methods can substantially outperform global methods (e.g., 39.8\% vs.\ 20.8\% precision at top-5\% in one controlled experiment). Under global contamination, the advantage disappears or reverses.

\textbf{Diagnostic heuristic.} Low agreement between a global detector (IF) and a contextual detector (peer-normalized IF) at a given threshold can indicate the presence of contextual structure in the data. This can inform method selection before committing to a full pipeline.

\subsection{KYC Integration in Practice}

\begin{itemize}
\item \textbf{Customer 360 Profiles}: Dynamic integration of KYC data, transaction history, external datasets for trend/anomaly identification and peer group comparison.
\item \textbf{HSBC Dynamic Risk Assessment}: ML models on transaction patterns + network behaviors + KYC data. Result: \textbf{60\% alert reduction, 2--4$\times$ true positive increase}.
\item \textbf{AI risk rating}: Incorporating behavioral and transactional factors into dynamic customer risk scores, enabling peer group comparisons.
\end{itemize}

\subsection{Landscape Assessment}

\textbf{The honest summary.} Contextual anomaly detection is a relatively thin literature:
\begin{itemize}
\item The core lineage is Song (2007) $\to$ ROCOD (2016) $\to$ QCAD (2023) $\to$ ConQuest (2024) $\to$ CWAE (2025).
\item Most AML papers do global anomaly detection. Very few model $P(\mathbf{y}|\mathbf{c})$ explicitly.
\item The deep contextual methods (CWAE, CVAE, CMDN) are recent (2024--2025) and largely untested on AML-specific data. King et al.\ (2025) benchmark on general tabular datasets (Bank, Census, KDD), not AML transaction data.
\item Peer group analysis is widely used in industry AML (Bolton and Hand, 2001) but rarely published as formal ML research---it sits in the practitioner/vendor space.
\item The context feature selection problem (which features are context vs.\ behavioral?) is largely unsolved. ConQuest's genetic algorithm and King et al.'s bilevel optimization are the only principled approaches.
\item \textbf{Gap}: No published work combines CWAE/WAE-MMD with AML-specific evaluation on transaction monitoring data with realistic laundering typologies. Similarly, peer-normalized IF with systematic evaluation across injection strategies is understudied.
\end{itemize}

%==============================================================================
\section{Large-Scale Entity Resolution}
%==============================================================================

Entity resolution (ER)---also called record linkage, deduplication, or identity resolution---is the task of determining which records across (or within) datasets refer to the same real-world entity. At AML scale (100M+ clients across multiple systems), naive pairwise comparison is $O(n^2)$ and infeasible. This section covers blocking strategies, matching methods, and scalable architectures.

\subsection{The Blocking Problem}

\textbf{Blocking is 10$\times$ more important than matching.} Reducing candidates from $n^2$ to $0.01 \times n^2$ makes everything tractable. Modern blocking approaches:

\begin{center}
\small
\begin{tabular}{p{3cm}p{4cm}p{4.5cm}}
\toprule
\textbf{Approach} & \textbf{Mechanism} & \textbf{Scale} \\
\midrule
Standard blocking & Hash on key attributes (name initial + DOB) & Simple; misses typos/variants \\
Sorted neighborhood & Sort by blocking key; slide window & Better recall; still key-dependent \\
LSH (Locality Sensitive Hashing) & Hash similar items to same bucket with high probability & Scalable; tunable recall/precision \\
ANN (Approximate Nearest Neighbor) & Embedding $\to$ FAISS/HNSW index $\to$ nearest neighbors & State-of-the-art; semantic similarity \\
Neural LSH & Learned hash functions via deep networks & Up to 6$\times$ faster index, 2$\times$ faster query (DET-LSH) \\
\bottomrule
\end{tabular}
\end{center}

\textbf{Neural LSH for Entity Blocking} (Amazon, 2024, arXiv:2401.18064). Proposes training deep neural networks as hashing functions for complex metrics. Significantly outperforms traditional LSH on entity resolution benchmarks.

\textbf{BlockingPy} (2025, arXiv:2504.04266). Open-source package implementing ANN blocking via FAISS (exact KNN, LSH, HNSW). Designed for entity resolution at scale with imperfect data.

\textbf{Hashed Dynamic Blocking.} Applied to a 530M-row industrial dataset, detecting 68 billion candidate pairs in under three hours.

\subsection{Matching Methods}

\subsubsection{Probabilistic: Fellegi-Sunter Model}

The classical approach. For each record pair, compute match weights from field-level agreement/disagreement probabilities:
\begin{equation*}
w_j = \log_2 \frac{P(\text{agree}_j | \text{match})}{P(\text{agree}_j | \text{non-match})}
\end{equation*}
Aggregate weights across fields; threshold for match/non-match/clerical review.

\textbf{Splink} (UK Ministry of Justice, open-source). The leading implementation of Fellegi-Sunter at scale. Key capabilities:
\begin{itemize}
\item Links 100M+ records using Spark or AWS Athena backends; 1M records on laptop in $\sim$1 minute.
\item Unsupervised (no labeled training data): learns match/non-match weights via EM algorithm.
\item Won Civil Service Awards 2025 (Innovation) and OpenUK Awards 2025.
\item Used by Harvard Medical School on 8.1M records (American Journal of Epidemiology, 2025).
\item Supports fuzzy string comparisons (Jaro-Winkler, Levenshtein), date comparisons, geographic distance.
\end{itemize}

\subsubsection{Deep Learning: Transformer-Based Matching}

The paradigm shift since 2020: use pre-trained language models (BERT, SentenceBERT, DistilBERT, RoBERTa) to encode record pairs into embeddings and classify match/non-match.

\textbf{Ditto} (Li et al., 2020). Landmark paper: BERT-based entity matching achieving SOTA, beating previous benchmarks by up to 29\% F1. Key idea: serialize record pairs as text, fine-tune BERT as binary classifier. Introduced data augmentation for ER (attribute-level).

\textbf{SentenceBERT for ER.} SentenceBERT models rank first on most ER benchmarks because they natively produce embedding vectors suitable for similarity search. Used both for blocking (embed $\to$ ANN search) and matching.

\textbf{EMBA} (2024, EDBT). Multi-task learning of BERT for entity matching, improving over single-task fine-tuning.

\textbf{Ensemble transformers} (2024). Combining multiple transformer models for better entity matching through ensembling.

\textbf{CSGAT} (2025, \textit{Scientific Reports}). Contextual Semantics Graph Attention Network extracting contextual information at token and attribute levels to generate semantically fused embeddings for entity matching.

\subsubsection{LLM-Based Entity Resolution}

The newest frontier (2024--2025):

\begin{itemize}
\item \textbf{Peeters and Bizer (2024)}: GPT-4 outperforms fine-tuned pre-trained language models by \textbf{40--68\%} on unseen entity types in zero-shot settings. No training data required.
\item \textbf{GPT-4o}: More cost-effective alternative; slightly lower performance than GPT-4.
\item \textbf{In-context clustering ER} (2025, arXiv:2506.02509): Packs multiple records into sets for direct in-context clustering via LLM. Generalizes SOTA approaches with improved cost-effectiveness.
\item \textbf{Balancing efficiency and quality} (Springer 2025): Explores optimal prompting strategies for LLM-based ER on structured data. ``Selecting'' strategy identified as most cost-effective.
\item \textbf{Limitation}: Cost at 100M+ scale---each LLM call is expensive. LLMs best used for high-uncertainty pairs after initial blocking+probabilistic matching.
\end{itemize}

\subsection{Scalable Architectures and Tools}

\begin{center}
\small
\begin{tabular}{p{2.2cm}p{2cm}p{3cm}p{4.5cm}}
\toprule
\textbf{Tool} & \textbf{Scale} & \textbf{Backend} & \textbf{Key Feature} \\
\midrule
Splink & 100M+ & Spark, Athena, DuckDB & Probabilistic (Fellegi-Sunter); unsupervised \\
Zingg & 100M+ & Spark, Snowflake, Databricks & ML-based blocking + matching; active learning \\
Senzing & Billions & In-memory graph & Real-time ER; entity-centric; used by banks \\
AWS ER & 100M+ & AWS Glue ML & Managed service; integrates with Zingg \\
Neo4j & 100M+ & Graph database & Graph-based ER; ownership chains \\
\bottomrule
\end{tabular}
\end{center}

\textbf{Zingg} (open-source, AGPL v3.0). ML-based entity resolution for Spark, Snowflake, Databricks. Learns blocking model that indexes near-similar records---typical comparisons are 0.05--1\% of the problem space. Supports both probabilistic (fuzzy) and deterministic matching. Active learning: human labels a small sample, model generalizes.

\textbf{Senzing} (commercial). In-memory, entity-centric ER engine. Processes billions of records in real-time. Used in financial crime, AML, and government applications. Key differentiator: maintains a live entity graph that updates incrementally as new records arrive---no batch reprocessing.

\subsection{ER Performance Benchmarks}

\textbf{Leipzig benchmark datasets} (Database Group Leipzig): Standard ER benchmarks with ground truth, including product, bibliographic, and geographic datasets at various scales.

\textbf{Key results on benchmarks:}
\begin{itemize}
\item BERT/Ditto (2020): +29\% F1 over previous SOTA on structured ER benchmarks.
\item SentenceBERT: Ranks 1st on most datasets; never below 4th.
\item GPT-4 (zero-shot): +40--68\% over fine-tuned PLMs on unseen entity types.
\item Splink: Near-SOTA accuracy at 100M+ scale with zero training data.
\end{itemize}

\textbf{The emerging architecture for 100M+ AML ER:}
\begin{enumerate}
\item \textbf{Embedding}: SentenceBERT or domain-specific encoder $\to$ vector per record
\item \textbf{Blocking}: ANN index (FAISS/HNSW) $\to$ candidate pairs (0.01--1\% of $n^2$)
\item \textbf{Matching}: Probabilistic (Splink) or fine-tuned transformer
\item \textbf{Uncertain pairs}: LLM (GPT-4) for high-ambiguity cases
\item \textbf{Graph consolidation}: Neo4j/knowledge graph for entity-level view
\end{enumerate}

%==============================================================================
\section{Feature Engineering}
%==============================================================================

\textbf{Feature engineering matters more than model choice.} The gap between well-engineered XGBoost and poorly-engineered DL is larger than between model architectures.

\begin{center}
\small
\begin{tabular}{p{2.5cm}p{5cm}p{4cm}}
\toprule
\textbf{Category} & \textbf{Examples} & \textbf{Impact} \\
\midrule
Velocity & Txn count/amount in 1h, 6h, 24h, 7d, 30d windows; unique counterparties; frequency acceleration; time since last txn & AUC +5--8\% (Bahnsen et al., 2016) \\
Behavioral & Amount ratio to hist.\ mean; Z-score; new merchant/country; dormancy; day/time deviation & Most important category \\
Network/Graph & Degree; PageRank; clustering coeff.; fan-in/fan-out; cycle detection; community membership & F1 +7\% even without GNNs (Alarab et al., 2020) \\
Aggregation & Per-customer stats (mean, std, skew); per-merchant fraud rate; per-country-pair amounts; per-channel stats & Foundation for behavioral deviation \\
Temporal & Day of week; hour; weekend/holiday; days since account open; seasonal patterns & Laundering timing patterns \\
Risk Indicators & FATF country risk; customer risk rating; PEP; industry risk; historical SAR count & Contextual conditioning \\
\bottomrule
\end{tabular}
\end{center}

\textbf{Two-stage feature architecture.} Gaining popularity:
\begin{enumerate}
\item Unsupervised scoring (IF anomaly scores, PCA reconstruction errors, autoencoder loss, cluster distances)
\item Supervised classification on original features + Stage 1 scores
\end{enumerate}
Addresses label scarcity. Jullum et al.\ (2020): $\sim$3\% AUC improvement.

\subsection{Handling Multicollinearity in Feature-Rich AML Pipelines}

\textbf{Detection.}
\begin{itemize}
\item \textbf{Pearson / Spearman correlation matrix}: Drop one of any pair with $|r| > 0.90$--$0.95$.
\item \textbf{Variance Inflation Factor (VIF)}: VIF $> 5$--$10$ indicates problematic collinearity.
\item \textbf{Eigenvalue analysis}: Condition number $> 30$ signals severe multicollinearity.
\end{itemize}

\textbf{Mitigation strategies.}

\begin{center}
\small
\begin{tabular}{p{3cm}p{4.5cm}p{4.5cm}}
\toprule
\textbf{Method} & \textbf{How It Works} & \textbf{When to Use} \\
\midrule
Correlation filter & Drop one of each highly correlated pair ($|r|>0.9$) & Quick first pass; interpretable \\
VIF-based removal & Iteratively remove highest-VIF feature until all VIF $<5$ & Linear models, logistic regression \\
Variable clustering (VARCLUS) & Hierarchical clustering of features; keep one representative per cluster & Large feature sets; preserves coverage of feature space \\
PCA / kernel PCA & Project onto orthogonal components & When interpretability of individual features is not required \\
Mutual information (mRMR) & Select features maximizing relevance to target while minimizing redundancy among selected features & Supervised setting with labels; captures nonlinear dependencies \\
LASSO / Elastic Net & $L_1$ penalty drives correlated feature coefficients to zero & Embedded selection; regularization path shows stability \\
Recursive Feature Elimination (RFE) & Iteratively remove least-important features using model importance & Works with any model; computationally expensive \\
Attribution-based selection & Cluster features by explanation-value correlation; keep representative of each cluster & Post-hoc; directly addresses attribution dilution \\
\bottomrule
\end{tabular}
\end{center}

%==============================================================================
\section{Datasets and Evaluation}
%==============================================================================

\subsection{Public Datasets}

\begin{center}
\small
\begin{tabular}{p{2.8cm}lp{2.2cm}lp{3.5cm}}
\toprule
\textbf{Dataset} & \textbf{Type} & \textbf{Size} & \textbf{Imbal.} & \textbf{Notes} \\
\midrule
Elliptic & Real BTC & 203K txns & 2\% & Most-used; 49 timesteps; 166 features \\
Elliptic++ & Real BTC & 204K+ & 2\% & +address-level graph; subgraph labels \\
AMLSim (IBM) & Synthetic & Config. & Config. & Open-source; controllable patterns \\
PaySim & Synthetic & 6.3M txns & 0.13\% & Mobile money; network structure \\
IBM AML Synth. & Synthetic & 180M+ & Config. & Large-scale multi-bank \\
IEEE-CIS & Real CC & 590K & 3.5\% & Kaggle competition; rich features \\
European CC & Real CC & 284K & 0.17\% & PCA-anonymized \\
\bottomrule
\end{tabular}
\end{center}

\textbf{Synthetic data.} Altman et al.\ (2023, \textit{Nature Scientific Data} / NeurIPS 2023): large-scale synthetic AML benchmark. Market projected \$310--576M (2024) $\to$ \$1.8--16.7B by 2030. By 2025, 75\% of large banks use synthetic data for AI. Tools: CTGAN, Copula GAN, TVAE---96--99\% utility equivalence.

\subsection{Class Imbalance}

AML: 0.1--2\% positive rate. Approaches:
\begin{itemize}
\item \textbf{SMOTE}: Most widely used. RF+SMOTE: F1 from 0.72 $\to$ 0.87 (Al-Hashedi and Magalingam, 2021).
\item \textbf{SMOTE-GAN hybrids}: SMOTified-GAN, etc.
\item \textbf{Cost-sensitive learning}: Class weights, focal loss.
\item \textbf{GraphSMOTE}: Graph-aware oversampling.
\item \textbf{AMLNet} (2025): SMOTE + RandomUnderSampler + multi-agent framework.
\end{itemize}

\subsection{Metrics}

AUC alone is insufficient for AML:
\begin{itemize}
\item \textbf{F1 (illicit)}: More discriminating on imbalanced data. Ranges 0.47 (RF) to 0.88 (PC-GNN).
\item \textbf{Precision@top-$k$}: How many top-$k$ alerts are true? Maps to analyst workload.
\item \textbf{Alert-level precision}: \% of alerts $\to$ SARs. Rules: $\sim$2\%; hybrid ML: $\sim$15\%.
\item \textbf{FP reduction}: Primary industry KPI. ML consistently 40--60\% over rules.
\item \textbf{PSI}: Monitors concept drift; triggers retraining.
\end{itemize}

\subsection{Concept Drift}

Launderers adapt $\Rightarrow$ model degradation:
\begin{itemize}
\item Deprez et al.\ (2025): Continual graph learning review---replay, regularization, architecture strategies.
\item Cross-regional continual learning with knowledge transfer (TKDE 2024).
\item PSI-based drift detection standard in production.
\item Traditional models suffer catastrophic forgetting on fine-tuning.
\end{itemize}

%==============================================================================
\section{Emerging Paradigms}
%==============================================================================

\subsection{Federated Learning}

Banks can't share data but laundering spans institutions:
\begin{itemize}
\item \textbf{FAML}: Multi-bank collaborative training without data sharing. Up to \textbf{30\% higher accuracy} vs.\ isolated.
\item \textbf{FedGraph-VASP} (2026): Federated graph learning + post-quantum crypto for cross-institutional AML.
\item One implementation on PaySim: 0.999 accuracy, 0.1\% false alarm rate.
\item \textbf{FHE + GNNs}: Fully homomorphic encryption for privacy-preserving graph ML.
\item \textbf{IBM Research}: Privacy-preserving federated learning for financial crime.
\item \textbf{Federated Graph Anomaly Detection via Disentangled Representations} (TheWebConf 2025).
\end{itemize}

\subsection{Explainability}

Regulatory non-negotiable (FinCEN, FCA, EBA, FATF):
\begin{itemize}
\item \textbf{SHAP/TreeSHAP}: Widely used. Jullum et al.\ (2020) reported analysts 30\% faster with SHAP explanations. However, SHAP values can be unstable with correlated features (common in AML), and computational cost grows with feature count. Should be used alongside---not as a replacement for---domain expertise.
\item \textbf{LIME}: An alternative to SHAP for local explanations, though less stable on correlated tabular features. Useful for quick per-alert explanations but sensitive to perturbation strategy.
\item \textbf{SEFraud} (KDD 2024): Self-explainable fraud via interpretative mask learning.
\item Counterfactual explanations (``this alert would not have fired if...'') demanded by regulators.
\item Global surrogates (decision trees/GAMs approximating black-box) for regulatory documentation.
\item EU AML Directive (2020), FinCEN guidance (2021): mandate ``understandable'' automated detection.
\end{itemize}

\subsection{Agentic AI}

Most recent frontier (2025--2026):
\begin{itemize}
\item \textbf{Nasdaq Verafin}: Agentic AI Workforce automating low-value, high-volume compliance. Up to 50\% AML workload automation.
\item \textbf{Co-Investigator AI} (2025): Agentic SAR generation---faster and more accurate.
\item McKinsey (2025): Agentic AI automates EDD reviews (adverse media, historic activity, escalation).
\item Agent-based KYC: onboarding from weeks $\to$ hours.
\item Human-in-the-loop: AI handles lower-risk; humans retain high-level decisions.
\end{itemize}

\subsection{Adversarial Robustness}

\begin{itemize}
\item CARE-GNN: Anti-camouflage (fraudsters connecting to benign nodes).
\item RoBCtrl (2025): RL-based attacks on GNN bot detectors---highlights deployed system vulnerability.
\item Adversarially robust temporal graph contrastive learning (2025).
\item Graph adversarial training (2024) for GNN robustness.
\end{itemize}

%==============================================================================
\section{Decision Framework}
%==============================================================================

\begin{center}
\small
\begin{tabular}{p{3cm}llp{4.5cm}}
\toprule
\textbf{Scenario} & \textbf{Data} & \textbf{Labels} & \textbf{Recommended Approach} \\
\midrule
Tabular, labeled & Tabular & Yes & XGBoost/LightGBM + post-hoc explanations \\
Tabular, no labels & Tabular & No & IF + two-stage \\
Graph, labeled & Network & Yes & GraphSAGE/GAT + temporal \\
Graph, no labels & Network & No & GAE or self-supervised GNN \\
Multi-institution & Any & Mixed & Federated learning (FAML) \\
Real-time & Stream & Any & TGN / streaming GNN \\
Regulatory focus & Any & Any & Hybrid rule+ML + explainability layer \\
Few labels & Any & Few & Self-supervised + fine-tuning \\
\bottomrule
\end{tabular}
\end{center}

%==============================================================================
\section{Industry Adoption}
%==============================================================================

\begin{center}
\small
\begin{tabular}{llll}
\toprule
\textbf{Organization} & \textbf{Method} & \textbf{Scale} & \textbf{Year} \\
\midrule
HSBC & Dynamic Risk Assessment (ML+KYC) & Banking & 2024 \\
Ant Group & Heterogeneous GNN & Billions of nodes & 2020+ \\
eBay & xFraud (Het-GNN) & Billions of txns & 2022 \\
PayPal & GNN + streaming & Real-time & 2021+ \\
Feedzai & Graph-enhanced ML & Banking AML & 2022+ \\
Elliptic & AI on Bitcoin & 200M txns & 2024 \\
Nasdaq Verafin & Agentic AI & Banking AML & 2025 \\
\bottomrule
\end{tabular}
\end{center}

\textbf{SAS/KPMG/ACAMS 2025 survey} (850 ACAMS members): 18\% have AI/ML in production; 18\% piloting; 25\% plan in 12--18 months; 40\% no plans. Top priority: FP reduction (38\%). ML top tech choice (58\%). Top barrier: lack of regulatory imperative (37\%, overtook budget). Regulator enthusiasm declining (51\% say regulators encourage, down 15pts from 2021).

%==============================================================================
\section{Supervised Learning with STR-Labeled Data}
%==============================================================================

When dispositioned STR (Suspicious Transaction Report) data is available, supervised learning becomes viable and is generally the highest-performing approach.

\subsection{Training Data Construction}

\begin{itemize}
\item \textbf{Label source}: Filed STRs = positive class; reviewed-and-cleared alerts = negative class. Raw non-alerted transactions are \emph{not} true negatives (unknown positives problem).
\item \textbf{Recommended window}: 6--24 months of dispositioned cases. Too short = insufficient positive examples; too long = concept drift contaminates the training set.
\item \textbf{Label quality}: STR filing decisions vary by analyst. Consensus labels (multiple reviewers) or adjudicated labels (senior review) improve training signal. Label noise is the single largest risk to supervised AML models.
\item \textbf{Class imbalance}: STR filing rates are typically 2--5\% of alerts, which are themselves $<$1\% of transactions. Effective mitigation: SMOTE, cost-sensitive learning ($w_{\text{pos}} = N_{\text{neg}} / N_{\text{pos}}$), or focal loss. GraphSMOTE extends oversampling to graph-structured data.
\end{itemize}

\subsection{Supervised Model Performance}

Gradient-boosted trees dominate supervised AML:
\begin{itemize}
\item XGBoost on tabular AML features: \textbf{AUROC 0.961} reported on production-scale labeled data (Jullum et al., 2020; DNB Bank study).
\item LightGBM: comparable AUROC with 2--3$\times$ faster training; preferred at scale.
\item Random Forest: strong baseline (AUROC $\sim$0.94) but less competitive than boosting.
\item Deep learning (DNN, LSTM): marginal gains over XGBoost on tabular data; justified only when sequential or graph structure is exploited.
\end{itemize}

\textbf{False positive reduction} is the primary value proposition of supervised models over rule-based systems. Literature consistently reports \textbf{60--80\% FP reduction} while maintaining or improving recall:
\begin{itemize}
\item Deloitte/SAS (2024): 60\% FP reduction with supervised ML overlay on existing rules.
\item Guidehouse (2023): ML models reduce alert volumes by 70\% with no loss in STR filing rate.
\item FinCEN (June 2024 proposed rule): explicitly endorses ML for ``greater precision in assessing customer risk'' and ``reducing false positives.''
\end{itemize}

\subsection{Semi-Supervised Extensions}

When labeled data is limited (common in early deployment):
\begin{itemize}
\item \textbf{Two-stage pipeline}: (1) unsupervised anomaly scores (IF, autoencoder, LOF) as features, (2) supervised classifier on the enriched feature set. This is the most common production architecture.
\item \textbf{Self-training / pseudo-labeling}: Train initial model on labeled subset, predict on unlabeled, add high-confidence predictions to training set. Iterate.
\item \textbf{Positive-unlabeled (PU) learning}: Treats non-STR transactions as unlabeled (not negative). Appropriate given the unknown-positives problem. Methods: Elkan \& Noto (2008), unbiased PU risk estimator (Kiryo et al., 2017).
\item \textbf{Contrastive pre-training}: Self-supervised on full transaction corpus, then fine-tune on STR labels. FLAG (KDD 2025), TREASURE (2025) demonstrate this for financial data.
\end{itemize}

%==============================================================================
\section{Complex Behavior Detection}
%==============================================================================

FINTRAC and international regulators increasingly focus on \emph{efficacy}---the ability to detect complex, multi-step laundering typologies, not just individual suspicious transactions.

\subsection{Laundering Typologies}

\begin{itemize}
\item \textbf{Structuring (smurfing)}: Breaking large amounts into sub-threshold deposits across accounts/branches. Rule-based detection (aggregation thresholds) is well-established but easily evaded by varying amounts.
\item \textbf{Layering}: Moving funds through multiple intermediary accounts/entities to obscure the trail. Requires multi-hop graph analysis.
\item \textbf{Mule networks}: Recruiting third-party accounts to receive and forward funds. Detected via fan-in/fan-out patterns and behavioral anomalies.
\item \textbf{Trade-based laundering (TBML)}: Over/under-invoicing in international trade. Requires trade data integration and price anomaly detection.
\item \textbf{Nested/shell structures}: Using corporate hierarchies to obscure beneficial ownership. Entity resolution is critical here.
\end{itemize}

\subsection{Detection Approaches by Typology}

\begin{center}
\small
\begin{tabular}{p{2.5cm}p{4cm}p{5cm}}
\toprule
\textbf{Typology} & \textbf{Best Approach} & \textbf{Key Reference} \\
\midrule
Structuring & Temporal aggregation + supervised & XGBoost with rolling-window features \\
Smurfing & Subgraph mining & GARG-AML (2025): second-order neighborhood scoring; SMoTeF: temporal-constrained flow \\
Layering & Multi-hop GNN / path analysis & FlowScope (AAAI 2020): dense subgraph on flow networks \\
Mule networks & Community detection + fan patterns & Heterogeneous GNN with role-aware embeddings \\
TBML & Price anomaly + cross-dataset & Trade price deviation models + customs data \\
\bottomrule
\end{tabular}
\end{center}

\subsection{Multi-Typology Supervised Models}

With STR labels, a single model can learn multiple typologies simultaneously:
\begin{itemize}
\item \textbf{Multi-label classification}: Each STR is tagged with typology code(s). Train XGBoost/LightGBM with multi-label targets. Enables typology-specific explanations.
\item \textbf{Hybrid LSTM-GraphSAGE}: Captures both sequential transaction patterns (LSTM) and network structure (GraphSAGE). Reported \textbf{95.4\% accuracy} on simulated multi-typology data.
\item \textbf{Hierarchical detection}: Level 1 = anomaly score (unsupervised), Level 2 = typology classifier (supervised), Level 3 = case-level aggregation. Each level adds specificity.
\end{itemize}

%==============================================================================
\section{FINTRAC Regulatory Alignment}
%==============================================================================

FINTRAC's evolving stance creates both opportunity and obligation for ML-based AML systems:
\begin{itemize}
\item \textbf{Efficacy mandate}: New violation for failure to ensure compliance programs are ``reasonably designed, risk-based and effective'' (2025). Shifts focus from checkbox compliance to demonstrable detection capability.
\item \textbf{Enforcement escalation}: 20 AMPs in 2025 alone, exceeding \$200M total---the largest in FINTRAC history. Inadequate transaction monitoring is a primary target.
\item \textbf{FinCEN endorsement of ML}: June 2024 proposed rule explicitly endorses AI/ML for ``greater precision in assessing customer risk'' and ``reducing false positives.'' Canadian regulators are expected to follow.
\item \textbf{OSFI E-23 Model Risk Management} (effective May 2027): Will require model validation, ongoing monitoring, and documentation for all models used in risk management---including AML. Architecture must be designed with validation in mind from day one.
\item \textbf{Explainability requirement}: Both FINTRAC and OSFI require that automated decisions be ``understandable.'' Post-hoc explanation methods (e.g., SHAP, LIME, surrogate models) can help satisfy this, though no single method is universally sufficient. Black-box models (deep GNNs, LLMs) require additional justification.
\end{itemize}

\bigskip
\noindent{\small\itshape Note: Some citation details (exact page/volume numbers) should be verified against original sources before formal publication. Performance metrics reflect literature-reported values under specific experimental configurations.}

\end{document}
